\documentclass[11pt]{article}
\usepackage[margin=.86in]{geometry}
\usepackage[T1]{fontenc}
\usepackage{lmodern,amsmath,amssymb,amsthm,booktabs,microtype,xcolor,hyperref}
\usepackage{fancyhdr}
\hypersetup{colorlinks=true,linkcolor=blue!40!black,urlcolor=blue!45!black,
pdftitle={Testing the quasi-RH moment transfer},pdfauthor={Research continuation for Jude Gomila}}
\pagestyle{fancy}\fancyhf{}\fancyhead[L]{\small TWIN LADDER / ARITHMETIC TRANSFER}
\fancyhead[R]{\small October 6, 2026}\fancyfoot[C]{\thepage}
\setlength{\headheight}{14pt}
\newtheorem{theorem}{Theorem}
\newtheorem{proposition}{Proposition}
\newtheorem{corollary}{Corollary}
\newcommand{\dd}{\,d}
\newcommand{\Ree}{\operatorname{Re}}
\newcommand{\E}{\mathcal E}
\newcommand{\D}{\mathcal D}
\newcommand{\Z}{\mathbb Z}
\newcommand{\Q}{\mathbb Q}
\title{Testing the quasi-RH moment transfer\\[4pt]
\large Exact arithmetic identities and the missing exponent saving}
\author{Research continuation for Jude Gomila}
\date{October 6, 2026}
\begin{document}
\maketitle
\begin{abstract}
We attempt the proposed transfer from the posted quasi-Riemann-hypothesis
moment estimates to the existing Vaughan Type-II packet. The attempt does
not produce an improved zero-free region or a proof of RH. We derive the
exact meromorphic generating function of the Type-II coefficients: every
nontrivial zero retains its full logarithmic-derivative residue. We check
the rational-to-Eisenstein coefficient lift, isolate the principal-row
restriction in the fourth moment, and calculate the loss in the available
inverse-moment amplification. The actual sixth-power amplifier gives an
individual inverse-sum exponent $11/12+\varepsilon$, weaker than the imported
$7/8$ conclusion. The latter yields the endpoint packet-energy bound
$\E_{\rm II}(T)=O(e^{3T/4})$, but no smaller exponent is established here.
Exact finite coefficient and rational-algebra checks accompany the proofs.
This is a working transfer audit, not an independently refereed result or
a claim of novelty.
\end{abstract}

\section{Inputs, scope, and the actual target}
The external input, denoted \textbf{QS}, is the posted assertion that zeta,
all Dirichlet $L$-functions, and the specified finite-order Hecke
$L$-functions over $F=\Q(\sqrt{-3})$ are zero-free in $\Ree s>7/8$,
apart from principal poles \cite{openai}. The source examined here is pinned
to commit \texttt{adc7f1241b42} (full hash in the source index).
The proof entry point and comparator configuration were inspected in the
preceding audit; the complete Lean dependency tree has not been rebuilt.
Every assertion below that uses QS is explicitly an implication from this
external input. The finite checks in this bundle do not verify QS.

The internal inputs are the October 1 Vaughan note v9 and its retained
v8 growth argument \cite{v9,v8}. Put
\[
 u=24\mathbf1_{[-1/48,1/48]},\quad
 g=\kappa^{-1/2}(-\partial_x^2+1/4)^2u^{*12},\quad
 k=g*\widetilde g,
\]
where
\[
 \kappa=\frac{11431397901559182284554757008547201}
 {2872446304320552960000}.
\]
This real even packet satisfies $\operatorname{supp}k\subset[-1/2,1/2]$,
$k(0)=1$, and
\begin{equation}\label{eq:K}
 K(w)=\int_{\mathbb R}k(x)e^{wx}\dd x
 =\frac{(1/4-w^2)^4}{\kappa}
 \left(\frac{\sinh(w/48)}{w/48}\right)^{24}.
\end{equation}
In particular $K$ has a fourth-order zero at $w=1/2$, is nonzero for
$0<\Ree w<1/2$, and has vertical decay $O((1+|\Im w|)^{-16})$ on
$|\Ree w|\le1/2$.

The sums and energies to be bounded are
\begin{align}
 S(L)&=\sum_{n\ge2}\frac{\Lambda(n)}{\sqrt n}k(\log n-L),\label{eq:S}\\
 \E(T)&=\int_T^{T+1}|S(L)|^2\dd L.\nonumber
\end{align}
At each window $L\in[T,T+1]$, fix $U=V=\lfloor e^{12T/25}\rfloor$ and let
\begin{align}
 a_U(r)&=\sum_{d\mid r,\ d>U}\mu(d),\nonumber\\
 D_T(L)&=\sum_{r>U,\ b>V}\frac{a_U(r)\Lambda(b)}{\sqrt{rb}}
                     k(\log(rb)-L),\label{eq:D}\\
 \E_{\rm II}(T)&=\int_T^{T+1}|D_T(L)|^2\dd L.\nonumber
\end{align}
Every displayed packet sum is finite. The v9 reduction proves
\begin{equation}\label{eq:error}
 \|S-D_T\|_{L^2[T,T+1]}
 \le 6\cdot10^{15}(T+2)e^{-T/50}\qquad(T\ge1).
\end{equation}
The remaining target is subexponential energy at all scales, not a finite
sample or an average over a different character family.

\section{Exact transfer: the Type-II term keeps every zero pole}
First hold arbitrary positive integer cutoffs $U,V$ fixed. Write
\[
 M_U(s)=\sum_{d\le U}\mu(d)d^{-s},\qquad
 L_V(s)=\sum_{b\le V}\Lambda(b)b^{-s}.
\]
These are entire Dirichlet polynomials; $L_V$ is not a Dirichlet
$L$-function. Let
\[
 d_{U,V}(n)=\sum_{rb=n}a_U(r)\Lambda(b)\mathbf1_{b>V},\qquad
 \D_{U,V}(s)=\sum_n d_{U,V}(n)n^{-s}.
\]

\begin{theorem}[Meromorphic Type-II identity]\label{thm:poles}
Initially for $\Ree s>1$, and then by meromorphic continuation,
\begin{align}
 \D_{U,V}(s)
 &=\bigl(1-\zeta(s)M_U(s)\bigr)
       \left(-\frac{\zeta'}{\zeta}(s)-L_V(s)\right)\label{eq:product}\\
 &=-\frac{\zeta'}{\zeta}(s)-L_V(s)
       +M_U(s)\zeta'(s)+\zeta(s)M_U(s)L_V(s).\label{eq:expanded}
\end{align}
At every nontrivial zero $\rho$ of multiplicity $m_\rho$, this function has
residue $-m_\rho$, independently of $U,V$.
\end{theorem}
\begin{proof}
Absolute convergence permits Dirichlet convolution in $\Ree s>1$:
\[
 \sum_r a_U(r)r^{-s}
 =\zeta(s)\sum_{d>U}\mu(d)d^{-s}
 =1-\zeta(s)M_U(s).
\]
Multiply by the series for the von Mangoldt tail to obtain
\eqref{eq:product}; expansion gives \eqref{eq:expanded}.
At a nontrivial zero, $M_U,L_V,\zeta,\zeta'$ are holomorphic. The only pole
in \eqref{eq:expanded} is the logarithmic derivative, of residue
$-m_\rho$. This also handles multiple zeros.
\end{proof}

This proves the exact coefficient version of Vaughan's identity:
\begin{equation}\label{eq:coeff}
 d_{U,V}=\Lambda-\Lambda_{\le V}
                 -\mu_{\le U}*\log+\mu_{\le U}*\Lambda_{\le V}*1.
\end{equation}
The final convolution cannot be discarded. The checker treats $\log p$ as
independent formal symbols and checks \eqref{eq:coeff} coefficient by
coefficient, with no floating-point logarithms.

Define the fixed-cutoff response by replacing $D_T$ in \eqref{eq:D} with
$D_{U,V}$. Its full Laplace transform, initially for $\Ree w>1/2$, is
\[
 \int_{\mathbb R}D_{U,V}(L)e^{-wL}\dd L
       =K(w)\D_{U,V}(1/2+w).
\]
Indeed the change of variables $x=\log n-L$ gives the factor $K(w)$.
The difference between this transform and
$-K(w)(\zeta'/\zeta)(1/2+w)$ is entire: the only possible additional
singularity is the pole of $\zeta$ and $\zeta'$ at $s=1$, which
\eqref{eq:K} removes. At $w=\rho-1/2$ with positive real part the residue is
\begin{equation}\label{eq:residue}
 -m_\rho K(\rho-1/2)\ne0.
\end{equation}
Thus neither the Vaughan split nor this pole-annihilating packet moves a
zero pole to half its real displacement. This does not rule out a deeper
arithmetic estimate for the sum.

\paragraph{Moving-cutoff qualification.}
The transform theorem holds with $U,V$ fixed. It is not a transform formula
for a response whose cutoffs vary with $L$. The moving-window energy
$\E_{\rm II}$ is related to $\E$ by \eqref{eq:error}, not by silently
substituting $U=U(L)$ into a Dirichlet series.

\section{What the imported zero bound rigorously supplies}
Let
\[
 \delta_* =\sup_\rho(\Ree\rho-1/2),\qquad
 \gamma_{\rm II}=\limsup_{T\to\infty}
                 \frac{\log(1+\E_{\rm II}(T))}{T}.
\]
The functional equation gives $0\le\delta_*\le1/2$; attainment of the
supremum is not assumed. The inherited explicit formula is
\begin{equation}\label{eq:explicit}
 S(L)=-\sum_\rho K(\rho-1/2)e^{(\rho-1/2)L}-A(L),\qquad L\ge1,
\end{equation}
where $A$ is the contribution of trivial zeros and is bounded, with
exponential decay. The zero sum is absolutely convergent because the
packet decays to order 16 vertically and $N(H)=O(H\log H)$.

\begin{corollary}[Endpoint energy bound from QS]\label{cor:energy}
Assuming QS,
\[
 \boxed{\E_{\rm II}(T)=O(e^{3T/4})\quad(T\ge1).}
\]
Moreover the inherited exact growth identity transfers to the moving
Type-II energy:
\begin{equation}\label{eq:gamma}
 \gamma_{\rm II}=2\delta_*\le3/4.
\end{equation}
No improvement to the number $3/4$ is established in this note.
\end{corollary}
\begin{proof}
Put $B=\sum_\rho|K(\rho-1/2)|<\infty$. Equation \eqref{eq:explicit} gives
$|S(L)|\le(B+\sup_{L\ge1}|A(L)|)e^{\delta_* L}$ for $L\ge1$.
Under QS and reflection, $\delta_*\le3/8$, so
$\E(T)\le C e^{3T/4}$. The triangle inequality and \eqref{eq:error} give
the same bound for $\E_{\rm II}$; the error is a decaying polynomial times
an exponential. The constant here is an existence bound; no numerical
evaluation of the infinite zero sum is asserted.

For clarity, the v8 growth identity $\gamma=2\delta_*$, with $\gamma$ defined
using $\E$, follows from the same upper bound and the exact Laplace formula
\[
 \int_{\log2-1/2}^{\infty}S(L)e^{-wL}\dd L
       =-K(w)\frac{\zeta'}{\zeta}(1/2+w).
\]
If $\gamma<2\delta_*$, choose $\gamma/2<b<\delta_*$. Cauchy--Schwarz on
each unit interval makes the integral holomorphic on $\Ree w>b$.
Some zero has $\Ree\rho-1/2>b$, and the nonzero residue
\eqref{eq:residue} contradicts that holomorphy. This is the existing v8
argument, included to expose its assumptions, not claimed as a new result.
Finally \eqref{eq:error} and
$x^2\le2y^2+2(x-y)^2$ in both directions show that $\E$ and
$\E_{\rm II}$ have the same nonnegative logarithmic growth exponent.
\end{proof}

\section{The rational and Eisenstein coefficients are different}
The new moment theorems are over ideals (represented by specified
generators) of $F=\Q(\sqrt{-3})$. Let $\chi=\chi_{-3}$ and let
\[
 \mu_F^{\rm norm}(n)=\sum_{N\mathfrak a=n}\mu_F(\mathfrak a).
\]
Quadratic factorization yields
\[
 \zeta_F(s)=\zeta(s)L(s,\chi),\qquad
 \mu_F^{\rm norm}=\mu*(\mu\chi),\qquad
 \boxed{\mu=\mu_F^{\rm norm}*\chi.}
\]
The last equality is coefficientwise, since
$(\mu\chi)*\chi=\mathbf1_{n=1}$.
\begin{center}
\begin{tabular}{lrr}
\toprule
Prime type & $\mu(p)$ & $\mu_F^{\rm norm}(p)$\\
\midrule
$p\equiv1\pmod3$ (split) & $-1$ & $-2$\\
$p\equiv2\pmod3$ (inert) & $-1$ & $0$\\
$p=3$ (ramified) & $-1$ & $-1$\\
\bottomrule
\end{tabular}
\end{center}
For example, $\mu_F^{\rm norm}(4)=-1$ and
$\mu_F^{\rm norm}(49)=1$, whereas $\mu(4)=\mu(49)=0$.
The corresponding von Mangoldt series satisfies
\[
 -\frac{\zeta_F'}{\zeta_F}(s)
 =\sum_n(1+\chi(n))\Lambda(n)n^{-s}.
\]
These are exact corrections, not reasons the lift is impossible. A
transfer must retain the additional convolution, its scale-dependent
cutoffs, and the resulting mixed terms. Substituting rational $\mu(n)$
for ideal M\"obius coefficients without them is invalid. Fixed excluded
Euler factors in the source require their own finite corrections as well.

\section{Attempt 1: isolate a row in the marked inverse moment}
The source's marked inverse lemma uses
\[
 M_u(Z^r)=Z^{-r/2}\sum_n\mu_F(n)\psi_u(n)W(Nn/Z^r),
\]
where $\psi_u$ is one common row character. A product $Q_u$ of normalized
prime-slot sums has total logarithmic length $z$. Its stated estimate is
\begin{equation}\label{eq:moment}
 \sum_{0<Nu\ll Z^m}|M_u(Z^r)Q_u|^2\ll Z^{m+\varepsilon},
\end{equation}
subject to fixed positive margins in
\begin{equation}\label{eq:constraints}
 r+2z<m,\qquad 2r+8z<3m.
\end{equation}
The prime-slot lists and coefficients are independent of the varying row.

Dropping all other terms in the positive sum is a valid way to select a
row, but pays the entire right side of \eqref{eq:moment}. With no slots,
$D=Z^r$, it only gives an unnormalized inverse-sum bound with exponent
$(m+r)/2>r$, no better than the trivial scale $D$.

Even granting the largest possible coherent prime-slot amplitude
$|Q_{u_0}|=Z^{z/2-o(1)}$, the resulting unnormalized upper-bound exponent is
\[
 \frac{r+m-z}{2}>r+\frac z2
\]
by the first inequality in \eqref{eq:constraints}. Smaller amplifier
amplitude worsens this estimate. This proves a limitation of this direct
term-extraction argument, not a limitation of every use of the moment
theorem.

\section{Attempt 2: use the actual sixth-power amplifier}
The source has an additional amplifier which must not be overlooked.
It replicates a sixth-power-free row $u$ by $ua^6$, preserving the sextic
character away from the declared coprimality zeros. Its exact column
identity includes the divisors of $\operatorname{rad}a$; simply erasing
these zeros would be incorrect. The resulting source estimate is
\begin{equation}\label{eq:sixth}
 \sum_{Nu\asymp U_0}|M_u(D)|^2
 \ll\max\{U_0,U_0^{1/6}D^{5(1+c)/6}\}(U_0D)^\varepsilon
\end{equation}
for any fixed $c>0$, over the stipulated sixth-power-free rows.

For a fixed row take bounded $U_0$. On choosing $c$ and the small losses
in terms of a prescribed $\varepsilon$, \eqref{eq:sixth} gives
\begin{equation}\label{eq:single}
 \left|\sum_n\mu_F(n)\psi_{u_0}(n)W(Nn/D)\right|
           \ll_{\varepsilon,u_0,W}D^{11/12+\varepsilon}.
\end{equation}
The exponent is $1/2+(5/6)/2=11/12$. Thus this is a real individual-row
saving, but it does not improve the already imported $7/8$ boundary.
The source itself uses this amplifier as one ingredient in a more
elaborate common-signal argument; \eqref{eq:single} is not presented as
a description of its entire proof.

\paragraph{An optimistic exponent worksheet.}
As a further stress test, grant a combination of the prime marks and
sixth-power copies with no extra mask or transfer cost. We do not claim
this combined estimate is a theorem. Normalize $r=1$ and allow the
closure of \eqref{eq:constraints}. The minimal admissible row exponent is
\[
 m(z)=\max\{1+2z,(2+8z)/3\}.
\]
Granting $Z^{m/6}$ copies and squared slot amplitude $Z^z$ gives the
optimistic inverse-sum exponent
\begin{equation}\label{eq:worksheet}
 \frac{1+(5/6)m(z)-z}{2}
 =\begin{cases}
 11/12+z/3,&0\le z\le1/2,\\
 7/9+11z/18,&z\ge1/2.
 \end{cases}
\end{equation}
Its infimum is still $11/12$. Strict admissibility margins only increase
the worksheet cost. This is an algebraic test of one proposed transfer,
not a universal lower bound for character amplification. Replacing the
sixth power by a square or cube would require a different invariance
identity and a compatible raw moment theorem; it is not a parameter
change licensed by the cited estimate.

\section{Attempt 3: apply the fourth moment or halve the scale}
The source's fourth-moment lemma is stated for $3/4\le\kappa_0\le1$.
With positive prime-slot length $z$ it requires
\[
 n_1+n_2+6\kappa_0z\le M,
 \qquad \beta_*\le(1+\kappa_0)/2\quad(\kappa_0<1),
\]
where $\beta_*$ is a supremum over the Hecke family, not just zeta.
The row family excludes principal inducing characters when $z=0$, and
excludes a specified finite character group $\Theta$ when $z>0$.
Consequently the principal rational row cannot simply be inserted.

This exclusion has analytic content. A principal normalized plain sum
at scale $D$ is generally of size $D^{1/2}$; two such sums have squared
product of scale $D_1D_2$. An unrestricted $Z^{M+\varepsilon}$ bound
would fail when the product of their lengths exceeds the row range.
The paper treats exceptional transformed rows by matched centering.
Our filter also removes a main density term, but no identity has been
proved that turns the full filtered M\"obius--von-Mangoldt Type-II form
into that permitted centered comparison.

QS reaches $\kappa_0=3/4$, the lower endpoint of the stated range.
Furthermore the source's low probe bound is $Z^{3/16+\varepsilon}$ and
its Mellin signal has affine exponent $C(s)=s-11/16$, so
$C(7/8)=3/16$. Reapplying the same bounds unchanged cannot improve the
boundary. A different parameter geometry or a better probe estimate is
not ruled out, but is not proved here.

\paragraph{Why two half-sized factors do not contract the exponent.}
Even in an optimistic separated model, grant both centered factors their
individual cancellation exponent $\beta$ and set $R=Z^u$, $B=Z^{1-u}$.
The normalized product has size
\[
 R^{\beta-1/2}B^{\beta-1/2}
       =Z^{\beta-1/2}.
\]
At $u=1/2$ each factor uses half the logarithmic scale, but there are two
factors. The squared-energy rate is $2\beta-1$, giving $3/4$ at
$\beta=7/8$. This calculation is an exponent diagnostic, not a claim
that these separated bounds already control every cut-off coefficient
in \eqref{eq:D}. Such cutoffs and the extra norm-lift convolution can
introduce further costs. The calculation supplies no saving beyond QS.

\section{The surviving arithmetic obligation and bridge interpretation}
All cross terms in the actual energy are
\begin{align*}
 \E_{\rm II}(T)
 =\sum_{r_1,b_1,r_2,b_2>U}
 \frac{a_U(r_1)a_U(r_2)\Lambda(b_1)\Lambda(b_2)}
 {\sqrt{r_1b_1r_2b_2}}\,
 \int_T^{T+1}k(\log(r_1b_1)-L)k(\log(r_2b_2)-L)\dd L.
\end{align*}
The packet support makes the sum finite. Its Gram kernel is positive
semidefinite, but the arithmetic coefficients have signs. Keeping those
signs is essential; positive averaging over unrelated rows does not
identify this quadratic form.

The concrete missing transfer is a centered, principal-sensitive estimate
for this quadratic form, or a faithful common-signal replacement, with a
strict saving relative to the possible rightmost-zero exponent. For
example, a proved arithmetic recurrence implying
\[
 \gamma_{\rm II}\le q\gamma_{\rm II},\qquad 0<q<1,
\]
would imply RH by \eqref{eq:gamma}. Neither that recurrence nor even a
smaller fixed bound $\gamma_{\rm II}<3/4$ has been obtained here. This
criterion was already proposed in v8; the present work tests specific
ways of trying to derive it from the newly posted estimates.

For the bridge $X(w)=\xi(1/2+w)$ and $R_a(w)=X(w-a)/X(w+a)$,
QS and the paired canonical product imply that $R_a$ is inner on
$\Ree w>0$ for $a\ge3/8$. To see the bound, group zeros into conjugate
and reflection pairs. For horizontal displacement $b$ and height $t$,
their ratio factors have nonnegative real Blaschke parameters $a-b$ and
$a+b$ whenever $a\ge|b|$. Each factor is contractive; locally convergent
paired products preserve this property. Boundary common factors are
filled by continuation. This is the standard shifted-$\xi$ framework
\cite{suzuki} applied to the imported smaller strip.

The existing twistor angular operator fixes the constant vector, so it
retains this same scalar arithmetic requirement. The new safe width is
useful, but the fixed-width reconstruction does not supply the signed
estimate at widths tending to zero. No new twistor positivity mechanism
has been constructed in this attempt.

\section{Verification, provenance, and status}
Running \texttt{python3 verify.py} uses only the Python standard library
and writes \texttt{verification.json}. It checks:
\begin{itemize}
\item 40,960 exact formal-log coefficients of \eqref{eq:coeff}, across
five cutoff pairs and all $1\le n\le8192$; omission of the final
convolution is deliberately rejected when it contributes.
\item The norm-lift inverse and recovery identities through $n=8192$,
including split, inert, and ramified prime controls.
\item The local residue $-m$ for six zero multiplicities using rational
polynomial models with nonconstant analytic cofactors.
\item Exact rational consistency of the amplification worksheet and
balanced-product exponents. The continuum worksheet claim is proved
by the two formulas in \eqref{eq:worksheet}, not inferred from a grid.
\end{itemize}
These checks validate finite identities and implementation consistency.
They do not verify the external analytic moment theorems, the complete
Lean build, or an all-scale bound stronger than the stated implication
from QS. The preceding source notes' interval-energy calculations were
read but not rerun; none is used as evidence for a new asymptotic bound.

The accompanying source index pins the external paper and records the
input bundle hash. No repository or deployed application was changed.
The successful outputs of this attempt are an explicit transfer map,
the pole-preservation proof, and a quantified account of the attempted
amplifier. \textbf{The exponent contraction and RH remain unproved.}

\begin{thebibliography}{9}\small
\bibitem{openai} OpenAI, \emph{The Quasi-Riemann Hypothesis: A Zero-Free
Half-Plane $\Re(s)>7/8$}, September 30, 2026. Pinned source in the
\texttt{openai/math} repository, commit cited in Section 1. Relevant
labels: \texttt{lem:marked}, \texttt{lem:inverse-amplification},
\texttt{lem:plain}, \texttt{prop:probe-low},
\texttt{lem:continuation-criterion}.
\url{https://github.com/openai/math/blob/main/lean/docs/003.md}.
\bibitem{v9} \emph{Vaughan reduction of the filtered prime energy},
research continuation for Jude Gomila, October 1, 2026.
User research bundle \texttt{rh\_vaughan\_v9.zip}; working derivation,
not independently refereed.
\bibitem{v8} \emph{The global-bound closure audit: Exact growth, a Laplace
obstruction, and a weaker sufficient target}, October 1, 2026.
Retained as \texttt{previous\_global\_closure\_v8.tex} in the v9 bundle.
\bibitem{suzuki} M. Suzuki, \emph{A canonical system of differential
equations arising from the Riemann zeta-function},
\url{https://arxiv.org/abs/1204.1827}.
\end{thebibliography}
\end{document}
